\section{USAMO 1997}
        \begin{enumerate}
        	\item (\textit{USAMO 1997}) Let $p_1,p_2,p_3,...$ be the prime numbers listed in increasing order, and let $x_0$ be a real number between $0$ and $1$. For positive integer $k$, define


 $x_{k}=\begin{cases}0&\text{ if }x_{k-1}=0\\ \left\{\frac{p_{k}}{x_{k-1}}\right\}&\text{ if }x_{k-1}\ne0\end{cases}$


 where $\{x\}$ denotes the fractional part of $x$. (The fractional part of $x$ is given by $x-\lfloor{x}\rfloor$ where $\lfloor{x}\rfloor$ is the greatest integer less than or equal to $x$.) Find, with proof, all $x_0$ satisfying $0<x_0<1$ for which the sequence $x_0,x_1,x_2,...$ eventually becomes $0$.


 

	\item (\textit{USAMO 1997}) Let $ABC$ be a triangle, and draw isosceles triangles $BCD, CAE, ABF$ externally to $ABC$, with $BC, CA, AB$ as their respective bases. Prove that the lines through $A,B,C$ perpendicular to the lines $\overleftrightarrow{EF},\overleftrightarrow{FD},\overleftrightarrow{DE}$, respectively, are concurrent.


 

	\item (\textit{USAMO 1997}) Prove that for any integer $n$, there exists a unique polynomial $Q$ with coefficients in $\{0,1,...,9\}$ such that $Q(-2)=Q(-5)=n$.


 

	\item (\textit{USAMO 1997}) To \textit{clip} a convex $n$-gon means to choose a pair of consecutive sides $AB, BC$ and to replace them by three segments $AM, MN,$ and $NC,$ where $M$ is the midpoint of $AB$ and $N$ is the midpoint of $BC$. In other words, one cuts off the triangle $MBN$ to obtain a convex $(n+1)$-gon. A regular hexagon $P_6$ of area $1$ is clipped to obtain a heptagon $P_7$. Then $P_7$ is clipped (in one of the seven possible ways) to obtain an octagon $P_8$, and so on. Prove that no matter how the clippings are done, the area of $P_n$ is greater than $\frac{1}{3}$, for all $n\ge6$.


 

	\item (\textit{USAMO 1997}) Prove that, for all positive real numbers $a, b, c,$


 $\frac{1}{a^3+b^3+abc}+\frac{1}{b^3+c^3+abc}+\frac{1}{a^3+c^3+abc} \le \frac{1}{abc}$.


 

	\item (\textit{USAMO 1997}) Suppose the sequence of nonnegative integers $a_1,a_2,...,a_{1997}$ satisfies 


 $a_i+a_j \le a_{i+j} \le a_i+a_j+1$


 for all $i, j \ge 1$ with $i+j \le 1997$. Show that there exists a real number $x$ such that $a_n=\lfloor{nx}\rfloor$ (the greatest integer $\le x$) for all $1 \le n \le 1997$.


 

\end{enumerate}